\section{Fractional energy compactness on moving measures}
\label{sec:height-compactness}

The planar limit identifies the support at order one.  To improve the
\emph{normalized} excess, we must retain the first-order normal displacement.
Two ingredients do this: a positive energy estimate obtained by testing with
an unnormalized height, and a finite-partition compactness argument for
functions whose underlying measures vary.  We state the resulting convergence
in terms of measures and moments; its formulation will also justify the
singular limiting equation in the next section.

\subsection{A local energy estimate}

For a measure $\sigma$, a bounded measurable set $U$, and a real measurable
function $f$, write
\begin{equation}
 \mathcal E_{\sigma,U}(f)
 =\iint_{U\times U}
       \frac{|f(x)-f(y)|^2}{|x-y|^{n+1}}\,d\sigma(x)d\sigma(y),
 \label{eq:height-energy-definition}
\end{equation}
when this nonnegative integral is finite.  Its integrand is assigned value
zero on the diagonal.  Positive-dimensional upper growth implies that
singletons have measure zero, so this convention has no effect.

The following statement uses the geometric and moment hypotheses supplied
by Proposition~\ref{prop:flat-limit}, together with explicit oscillation
and horizon conditions.  The small parameter is denoted by
$\delta_j$; the integer $K_j$ records the last available dyadic moment.

\begin{lemma}[Energy of a normalized affine height]
\label{lem:height-energy}
Let $1\le n<d$, $0<\gamma<1$, and let $\mu_j$ be Radon measures on $\R^d$
with $0\in\supp\mu_j$.  Suppose that, for fixed $G,c>0$,
\begin{align}
 \mu_j(B(x,r))&\le Gr^n
       &&(x\in\R^d, r>0),\notag\\
 \mu_j(B(0,r))&\ge cr^n
       &&(0<r\le\diam\supp\mu_j),
 \label{eq:height-growth}
\end{align}
and $\diam\supp\mu_j\to\infty$.
Let $u_j$ be real affine functions with
\begin{equation}
 u_j(x)-u_j(y)=\langle e_j,x-y\rangle,
 \qquad |e_j|\le1,\qquad |u_j(0)|\le W.
 \label{eq:height-affine-data}
\end{equation}
Assume $\delta_j>0$, $\delta_j\to0$, $K_j\to\infty$, and
\begin{equation}
 \frac{2^{-K_j}}{\delta_j}\longrightarrow0.
 \label{eq:height-last-radius}
\end{equation}
Suppose, with a fixed $M<\infty$, that
\begin{equation}
 \int_{B(0,2^k)}|u_j|^2\,d\mu_j
 \le M\delta_j^2 2^{(n+2+2\gamma)k}
 \qquad(0\le k\le K_j).
 \label{eq:height-dyadic-moments}
\end{equation}
Finally, let $A_j\to\infty$ and suppose that every compactly supported
Lipschitz function $\phi$ with
$\supp\phi\subset B(0,A_j)$ and $\int\phi\,d\mu_j=0$ satisfies
\begin{equation}
 |\mathcal R_{\mu_j}(\phi)|
       \le \delta_j^3\Lip(\phi).
 \label{eq:height-small-oscillation}
\end{equation}
Then $w_j=u_j/\delta_j$ has finite local energy, and for every
$H<\infty$ there is $C_H<\infty$, independent of $j$, such that
\begin{equation}
 \int_{B(0,H)}|w_j|^2\,d\mu_j
       +\mathcal E_{\mu_j,B(0,H)}(w_j)\le C_H.
 \label{eq:height-uniform-energy}
\end{equation}
The eventual bound depends only on $H,n,\gamma,G,c,M,W$; finitely many
initial terms may be included by enlarging $C_H$.
\end{lemma}

\begin{proof}
It suffices to consider $H\ge1$.  Fix a Lipschitz function $\chi$ with
$0\le\chi\le1$, $\Lip(\chi)\le2/H$, equal to one on $B(0,H)$,
and zero outside $B(0,3H/2)$.  Its support lies in $B(0,2H)$.
Choose the smallest dyadic radius $R\ge4H$, and put $D=B(0,R)$.
All estimates below concern sufficiently large $j$, so that $K_j\ge\log_2R$,
$\delta_j\le1$, $A_j>2H$, and the lower estimate in
\eqref{eq:height-growth} is available at radius $H$.  Constants may depend
on the fixed quantities displayed in the lemma and on $H,R$, but not on $j$.
We omit the index $j$ temporarily.

The moment bound at radius $R$ gives
\begin{equation}
 \int_D u^2\,d\mu\le C\delta^2,
 \qquad
 \int\chi^2\,d\mu\ge\mu(B(0,H))\ge cH^n>0.
 \label{eq:height-cutoff-mass}
\end{equation}
Define the weighted center and its centered test by
\begin{equation}
 b=\frac{\int\chi^2u\,d\mu}{\int\chi^2\,d\mu},
 \qquad U=u-b,\qquad \phi=\chi^2U.
 \label{eq:height-centered-test}
\end{equation}
Cauchy--Schwarz, followed by \eqref{eq:height-cutoff-mass}, proves
\begin{equation}
 |b|\le C\delta,\qquad
 \int_DU^2\,d\mu\le C\delta^2,\qquad
 \int|\phi|\,d\mu\le C\delta,
 \qquad \int|\phi U|\,d\mu\le C\delta^2.
 \label{eq:height-centered-bounds}
\end{equation}
In particular, these estimates precede any energy bound.  The function
$\phi$ has mean zero.  It also has a Lipschitz constant bounded independently
of $j$: on the support of $\chi$,
$|U(x)|\le W+2H+C$, and the product rule for Lipschitz functions gives the
bound.  The same bound across the complement follows because $\chi$ vanishes
there.  Consequently \eqref{eq:height-small-oscillation} implies
\begin{equation}
 |\langle e,\mathcal R_\mu(\phi)\rangle|\le C\delta^3.
 \label{eq:height-tested-pairing}
\end{equation}

The positive part of this pairing comes from the exact identity
\begin{align}
 &(U(x)-U(y))(\chi(x)^2U(x)-\chi(y)^2U(y))\notag\\
 &\hspace{12mm}=
     |\chi(x)U(x)-\chi(y)U(y)|^2
       -U(x)U(y)(\chi(x)-\chi(y))^2.
 \label{eq:height-cutoff-identity}
\end{align}
All terms in its integral over $D\times D$ are integrable before a
quantitative energy estimate is made.  Indeed, a Lipschitz difference
cancels one power of the distance in each factor.  For every $x\in D$,
upper growth and dyadic annuli imply
\begin{equation}
 \int_D |x-y|^{1-n}\,d\mu(y)\le C G R.
 \label{eq:height-near-row}
\end{equation}
The integrability of local affine and cutoff energies follows from this
bound and the finiteness of $\mu(D)$.  For the signed cutoff error, symmetry
and $2|U(x)U(y)|\le U(x)^2+U(y)^2$ give the sharper estimate
\begin{align}
 &\iint_{D\times D}
   |U(x)U(y)|\frac{|\chi(x)-\chi(y)|^2}{|x-y|^{n+1}}
                      \,d\mu(x)d\mu(y)\notag\\
 &\hspace{12mm}\le C G R\Lip(\chi)^2\int_D U^2\,d\mu
       \le C\delta^2.
 \label{eq:height-cutoff-error}
\end{align}
This is an estimate from the already known second moment; it assumes no
bound on the unknown fractional energy.

We next separate the local pairing from its exterior remainder.  Set
$q(x,y)=|x-y|^{-n-1}$.  The affine identity in
\eqref{eq:height-affine-data}, and the pairing formula in
Lemma~\ref{lem:pairings}, give
\begin{align}
 \langle e,\mathcal R_\mu(\phi)\rangle
   &=\frac12\iint_{D\times D}
        (U(x)-U(y))(\phi(x)-\phi(y))q(x,y)
                         \,d\mu(x)d\mu(y)+\mathrm{Ext},
 \label{eq:height-pairing-split}\\
 \mathrm{Ext}
   &=\int_D\phi(x)\int_{D^c}
       \bigl[U(x)q(x,y)
           -U(y)\{q(x,y)-q(0,y)\}\bigr]\,d\mu(y)d\mu(x).
 \label{eq:height-exterior-cancelled}
\end{align}
To obtain the second identity, the base-point subtraction initially also
contributes $-U(0)q(0,y)$.  That term is separately integrable on $D^c$
and disappears because $\int\phi\,d\mu=0$.  The difference multiplying
$U(y)$ in \eqref{eq:height-exterior-cancelled} is kept together.

For $|x|\le2H$ and $|y|\ge R$,
\begin{equation}
 q(x,y)\le C|y|^{-n-1},
 \qquad |q(x,y)-q(0,y)|\le CH|y|^{-n-2}.
 \label{eq:height-tail-kernel}
\end{equation}
The global affine bound $|U(y)|\le |y|+W+C$, together with upper
growth and \eqref{eq:height-tail-kernel}, makes the exterior expression
absolutely integrable.  We may therefore integrate
\eqref{eq:height-cutoff-identity} and use
\eqref{eq:height-pairing-split} and \eqref{eq:height-cutoff-error} to obtain
\begin{equation}
 \mathcal E_{\mu,D}(\chi U)
 \le 2|\langle e,\mathcal R_\mu(\phi)\rangle|
       +2|\mathrm{Ext}|+C\delta^2.
 \label{eq:height-energy-reduction}
\end{equation}
It remains to bound the exterior remainder at the scale of the known
second moment.  Put $T=2^K$.
Upper growth and dyadic annuli give
$\int_{D^c}|y|^{-n-1}\,d\mu(y)\le C/R$.
On a controlled shell $t\le |y|<2t$, where $t=2^k$ and $k<K$,
Cauchy--Schwarz and the moment at radius $2t$ give
\begin{equation}
 \int_{t\le|y|<2t}|U(y)|\,d\mu(y)
          \le C\delta t^{n+1+\gamma}
 \qquad(t\ge R).
 \label{eq:height-shell-first-moment}
\end{equation}
Here the contribution of $b$ is absorbed using $t\ge1$.
Since $\gamma<1$, summing the controlled shells gives
\begin{equation}
 \int_{R\le|y|<T}\frac{|U(y)|}{|y|^{n+2}}\,d\mu(y)
       \le C\delta\sum_{\substack{t\ge R\\t\text{ dyadic}}}t^{\gamma-1}
       \le C\delta.
 \label{eq:height-controlled-tail}
\end{equation}
Beyond the actual last controlled radius, use the global affine bound
$|U(y)|\le |y|+W+C$.  Upper growth then gives
\begin{equation}
 \int_{|y|\ge T}\frac{|U(y)|}{|y|^{n+2}}\,d\mu(y)
       \le C(T^{-1}+T^{-2})\le C T^{-1}.
 \label{eq:height-uncontrolled-tail}
\end{equation}
The half-open shells and the closed final exterior cover every point,
including their boundary spheres.  Combining
\eqref{eq:height-centered-bounds} with the kernel and tail bounds
\eqref{eq:height-tail-kernel}--\eqref{eq:height-uncontrolled-tail}
gives the quantitative estimate
\begin{equation}
 |\mathrm{Ext}|\le C(\delta^2+\delta T^{-1}).
 \label{eq:height-exterior-bound}
\end{equation}
Thus the remote part retains the factor $\delta$ supplied by the centered
test.  Replacing that factor by a constant would lose the required
normalization.

The reduction \eqref{eq:height-energy-reduction}, the tested-pairing bound
\eqref{eq:height-tested-pairing}, and \eqref{eq:height-exterior-bound}
now yield
\begin{equation}
 \mathcal E_{\mu,D}(\chi U)
       \le C(\delta^3+\delta^2+\delta T^{-1}).
 \label{eq:height-unnormalized-energy}
\end{equation}
On $B(0,H)$, $\chi=1$ and the centered difference $U(x)-U(y)$ is
$u(x)-u(y)$.  Restriction to this smaller product set and division by
$\delta^2>0$ give
\begin{equation}
 \mathcal E_{\mu,B(0,H)}(u/\delta)
       \le C\left(1+\delta+\frac{1}{\delta T}\right).
 \label{eq:height-normalized-bound}
\end{equation}
The last term is bounded by \eqref{eq:height-last-radius}.  The local
$L^2$ estimate follows directly from the moment bound at a fixed dyadic
radius containing $H$.

For each of the finitely many omitted indices, $u_j/\delta_j$ is an
ordinary Lipschitz function.  Upper growth gives its finite $L^2$ norm
on every fixed ball, and \eqref{eq:height-near-row} gives its finite local
energy.  Enlarging the bound to cover those indices completes the proof.
\end{proof}

In our application $\delta_j=2^{-j}$ and $K_j=2j+3$.  In particular,
$\delta_j^{-1}2^{-K_j}=2^{-j-3}\to0$.  The conclusion of
Lemma~\ref{lem:height-energy} is a \emph{bounded} normalized energy;
no vanishing of that energy has been asserted.

\subsection{Finite partitions on bounded projection regions}

Fix a linear $n$-plane $L$ through the origin, an isometry
$e:\R^n\to L$, and the orthogonal coordinate map $\pi=e^*$, so that
$e\pi$ is the orthogonal projection onto $L$.  Set
\begin{equation}
 Q_H=(-(H+1),H+1)^n,
 \qquad
 \Omega_H=\pi^{-1}(Q_H)\cap
 B\bigl(0,(\sqrt n+1)(H+1)\bigr),
 \qquad H=0,1,2,\ldots .
 \label{eq:height-projection-regions}
\end{equation}
These bounded open sets exhaust $\R^d$.  Their intersection with $L$ is
$e(Q_H)$: the closed box $e(\overline Q_H)$ lies strictly inside the
ambient bounding ball, because $|e(z)|\le\sqrt n(H+1)$ there.  Consequently
\begin{equation}
 \nu(\partial\Omega_H)=0
 \quad\text{for}\quad \nu=q\,e_\#\mathcal L^n,\quad q>0.
 \label{eq:height-region-null-boundary}
\end{equation}
Only the coordinate faces meet the plane on this boundary, and those faces
have zero $n$-dimensional volume.

\begin{lemma}[Compactness with moving measures]
\label{lem:moving-compactness}
Suppose $\mu_j$ are Radon measures with a common global upper $n$-growth
bound and a common lower bound $cr^n$ at every support point and every
admissible radius.  Suppose also that
$\diam\supp\mu_j\to\infty$ and
\begin{equation}
 \mu_j\longrightarrow\nu=q\,e_\#\mathcal L^n,
       \qquad q>0,
 \label{eq:height-flat-measure-limit}
\end{equation}
against continuous compactly supported tests.  Let
$w_j^1,\ldots,w_j^s$ be finitely many real measurable functions.  Assume
that each belongs to $L^2$ on bounded balls and has integrable local
fractional energy, and that for every $R<\infty$
\begin{equation}
 \sup_j\sum_{a=1}^s
 \left(\int_{B(0,R)}|w_j^a|^2\,d\mu_j
        +\mathcal E_{\mu_j,B(0,R)}(w_j^a)\right)<\infty.
 \label{eq:height-compactness-hypothesis}
\end{equation}
There are one subsequence and measurable functions
$v^1,\ldots,v^s\in L^2_{\mathrm{loc}}(L,\nu)$ such that, on that
subsequence, for every $H$ and every continuous function $\psi$ on the
compact set $\overline\Omega_H$,
\begin{align}
 \int_{\Omega_H}w_j^a\psi\,d\mu_j
       &\longrightarrow\int_{\Omega_H}v^a\psi\,d\nu,
 \label{eq:height-signed-moment-limit}\\
 \int_{\Omega_H}|w_j^a|^2\psi\,d\mu_j
       &\longrightarrow\int_{\Omega_H}|v^a|^2\psi\,d\nu.
 \label{eq:height-square-moment-limit}
\end{align}
The ordinary restrictions $\mu_j|_{\Omega_H}$ also converge weakly to
$\nu|_{\Omega_H}$ as finite measures on $\overline\Omega_H$.
In particular,
\begin{equation}
 \int_{\Omega_H}|w_j^a|^2\,d\mu_j
       \longrightarrow\int_{\Omega_H}|v^a|^2\,d\nu.
 \label{eq:height-norm-limit}
\end{equation}
The subsequence and the functions are common to all these localizations
and tests.
\end{lemma}

\begin{proof}
We give the construction, including the norm convergence.
On each bounded projection region, fractional energy controls the
variance inside small cells. A common subsequence of their finitely many
averages will recover both first and second moments, not just weak
convergence of the heights.
First, \eqref{eq:height-region-null-boundary} and compact support cutoffs
imply finite-measure weak convergence on each $\Omega_H$.  Indeed, multiply
by a compact continuous cutoff equal to one on $\overline\Omega_H$;
weak convergence of these finite measures passes to restriction to the
continuity set $\Omega_H$.  In particular their masses are bounded.

We will also use the consequence
\begin{equation}
 \sup\{\dist(x,L):x\in\supp\mu_j\cap B(0,R)\}
                       \longrightarrow0
 \quad(R<\infty).
 \label{eq:height-physical-tube}
\end{equation}
If this failed, some bounded sequence of support points would stay a
positive distance from $L$.  Passing to a subsequence, a fixed small ball
about its limit would be disjoint from $L$ and would contain a smaller
ball about each of those support points.  The lower growth bound would
give the smaller ball a fixed positive mass; its radius is admissible
for large $j$.  A continuous compact cutoff, equal to one on all the
smaller balls and supported away from $L$, would then contradict
\eqref{eq:height-flat-measure-limit}.  The cutoff is taken on a larger
ambient ball, so the smaller support balls retain their full mass even
if they cross a localization boundary.

Fix $H$, and partition the half-open box
$(-(H+1),H+1]^n$ into equal dyadic boxes of side length
$h=2(H+1)2^{-k}$.  For such a box $D$, set
\[
 C_D=\Omega_H\cap\pi^{-1}(D).
\]
These finitely many disjoint measurable cells cover $\Omega_H$.  Each
cell boundary has $\nu$-measure zero, by
\eqref{eq:height-region-null-boundary} and the nullity of coordinate
faces.  The same holds for intersections of cells from different
partition levels.  Hence
\begin{equation}
 \mu_j(C_D)\longrightarrow\nu(C_D)=q h^n,
 \qquad
 \mu_j(C_D\cap C_{D'})\longrightarrow\nu(C_D\cap C_{D'}).
 \label{eq:height-cell-mass-limits}
\end{equation}
For each fixed partition level, all its cell masses are therefore at
least $(q/2)h^n$ for sufficiently large $j$.  This uses positivity on
\emph{every} full-plane cell, rather than just nonzero total mass.
By \eqref{eq:height-physical-tube}, at that same fixed level the actual
support points in one cell have mutual distance at most
$(\sqrt n+2)h$ for sufficiently large $j$.  The entire projection cylinder
need not have bounded diameter; only pairs seen by the restricted measure
are used.

For $w_j=w_j^a$, let $m_{j,D}$ be its average on $C_D$ (zero if the
cell has zero mass), and let $P_k^jw_j$ be the corresponding step function.
The exact variance
identity gives
\begin{align}
 \int_{C_D}|w_j-m_{j,D}|^2\,d\mu_j
 &=\frac{1}{2\mu_j(C_D)}
    \iint_{C_D\times C_D}|w_j(x)-w_j(y)|^2\,d\mu_j(x)d\mu_j(y)
       \notag\\
 &\le C_n q^{-1}h
    \iint_{C_D\times C_D}
       \frac{|w_j(x)-w_j(y)|^2}{|x-y|^{n+1}}
                     \,d\mu_j(x)d\mu_j(y).
 \label{eq:height-cell-variance}
\end{align}
All averages are defined for these sufficiently large indices, and their
$L^2$ bounds follow from \eqref{eq:height-compactness-hypothesis}.
The sets $C_D\times C_D$ are disjoint subsets of
$\Omega_H\times\Omega_H$.  Summing
\eqref{eq:height-cell-variance} yields
\begin{equation}
 \int_{\Omega_H}|w_j-P_k^jw_j|^2\,d\mu_j
                 \le C_H h.
 \label{eq:height-partition-error}
\end{equation}
There is no factor counting the cells in this estimate.

For fixed $H,k,D,a$, Cauchy--Schwarz and the positive limiting cell mass
bound $m_{j,D}^a$.  A diagonal subsequence therefore makes all these
averages converge simultaneously: the collection of indices is
countable, including all $H,k$ and the finite set of coordinates $a$.
Write $m_D^a$ for their limits and $P_{H,k}^a$ for the corresponding step
function on the limiting plane box.

At fixed levels $k,l$, expand the squared difference of the two step
functions.  Convergence of the finitely many coefficients, cell masses,
and intersection masses in \eqref{eq:height-cell-mass-limits} gives
\[
 \int_{\Omega_H}|P_k^jw_j^a-P_l^jw_j^a|^2\,d\mu_j
       \longrightarrow
       \int_{\Omega_H}|P_{H,k}^a-P_{H,l}^a|^2\,d\nu.
\]
The triangle inequality and \eqref{eq:height-partition-error} show that
the last integral is at most $2C_H(h_k+h_l)$.  The step functions thus
converge in $L^2(\nu|_{\Omega_H})$ to some $v_H^a$.

The reverse triangle inequality also gives
\[
 \left|\|w_j^a\|_{L^2(\mu_j|_{\Omega_H})}
       -\|P_k^jw_j^a\|_{L^2(\mu_j|_{\Omega_H})}\right|
       \le (C_H h_k)^{1/2}.
\]
For fixed $k$, the step-function norm converges, since its square is the
finite sum $\sum_D(m_{j,D}^a)^2\mu_j(C_D)$.  First let $j\to\infty$,
then $k\to\infty$.  This proves
\eqref{eq:height-norm-limit} with $v_H^a$ in place of $v^a$.

For a continuous $\psi$ on $\overline\Omega_H$, the signed testing error
from replacing $w_j^a$ by $P_k^jw_j^a$ is at most
\[
 (C_Hh_k)^{1/2}\|\psi\|_\infty\mu_j(\Omega_H)^{1/2}.
\]
For fixed $k$, weak convergence on the continuity cells gives
$\int_{C_D}\psi\,d\mu_j\to\int_{C_D}\psi\,d\nu$.  Passing first in
$j$, then in $k$, proves \eqref{eq:height-signed-moment-limit} for
$v_H^a$.

The local limits agree on overlaps.  To see this without choosing new
subsequences, let $U=\Omega_H\cap\Omega_J$ and
$\eta(x)=\min\{1,\dist(x,U^c)\}$.  This bounded Lipschitz function is
positive exactly on $U$.  For a compactly supported Lipschitz $\psi$,
use $\eta\psi$ in the two signed-moment limits.  The source integrals
are identical, since $\eta\psi$ vanishes outside the overlap.  Their
limits imply
\[
 \int_U\eta(v_H^a-v_J^a)\psi\,d\nu=0.
\]
Compact Lipschitz tests are dense in $L^2(\nu|_U)$, so
$\eta(v_H^a-v_J^a)=0$ almost everywhere.  Since $\eta>0$ on $U$,
the local limits agree there.  Choosing measurable representatives and,
on the plane, using the first region containing a point now produces
one measurable $v^a$ agreeing with every $v_H^a$ almost everywhere.
Countability of the cover ensures this simultaneous agreement.

It remains to prove the weighted squared-moment conclusion.  Fix $H,a$,
and abbreviate $w_j=w_j^a$, $v=v^a$.  There are bounded compactly supported
ambient Lipschitz functions $g$ whose restrictions approximate $v$ in
$L^2(\nu|_{\Omega_H})$.  For completeness, simple functions are dense
in this space; regularity of finite plane volume approximates their sets
by compact and open sets, and distance cutoffs approximate their
indicators.  Extend the resulting plane Lipschitz functions by $\pi$,
and multiply by a fixed ambient compact cutoff equal to one on
$\overline\Omega_H$.

For each such fixed $g$, the norm limit, the signed limit tested against
$g$, and ordinary weak convergence tested against $g^2$ give
\begin{equation}
 \lim_j\int_{\Omega_H}|w_j-g|^2\,d\mu_j
       =\int_{\Omega_H}|v-g|^2\,d\nu.
 \label{eq:height-continuous-approximation}
\end{equation}
This identity involves $g$, an actual continuous function, on the moving
supports.  For any continuous weight $\psi$,
\begin{align*}
 \left|\int_{\Omega_H}\psi(w_j^2-g^2)\,d\mu_j\right|
 &\le \|\psi\|_\infty
     \|w_j-g\|_{L^2(\mu_j|_{\Omega_H})}\\
 &\hspace{8mm}\cdot
     \left(\|w_j\|_{L^2(\mu_j|_{\Omega_H})}
                  +\|g\|_{L^2(\mu_j|_{\Omega_H})}\right).
\end{align*}
The last factor remains bounded for an $L^2$-approximating family $g$,
after taking $j\to\infty$ for each fixed $g$.  Ordinary weak convergence
handles $\psi g^2$.  Equation~\eqref{eq:height-continuous-approximation}
and Cauchy--Schwarz for the limiting measure then allow $g\to v$ in
$L^2(\nu|_{\Omega_H})$.  This proves
\eqref{eq:height-square-moment-limit} in the stated order of limits.
\end{proof}

The lemma may equivalently be read as three finite-measure convergences
on each common compact closure:
\begin{equation}
 \mu_j|_{\Omega_H}\ \longrightarrow\ \nu|_{\Omega_H},\qquad
 w_j^a\mu_j|_{\Omega_H}\ \longrightarrow\ v^a\nu|_{\Omega_H},\qquad
 |w_j^a|^2\mu_j|_{\Omega_H}\ \longrightarrow\ |v^a|^2\nu|_{\Omega_H}.
 \label{eq:height-three-measures}
\end{equation}
The middle measures are signed; their total variations are uniformly
bounded by Cauchy--Schwarz.  These statements do not require evaluating
an arbitrary representative of $v^a$ against $\mu_j$.  Once the next
section identifies $v^a$ with an affine function, the same polarization
as in \eqref{eq:height-continuous-approximation} will give a genuine
squared-error limit for that fixed affine representative.
